% !TEX program = xelatex
% ============================================================
%  صفحة واحدة (One-Pager) Template
%  قالب تسويقي مكثف في صفحة واحدة: مشكلة، حل، ميزات، إحصائيات
%  Compile with: xelatex main.tex (run twice)
% ============================================================
%  Features:
%    - Company/product header with tagline
%    - Problem statement and solution panels (tcolorbox)
%    - Key features (3-4 bullets with icon placeholders)
%    - Metrics/stats row (3 big numbers)
%    - Target market section
%    - Contact/CTA footer panel
%    - Visually dense but clean single-page layout
% ============================================================

\documentclass[11pt,a4paper]{article}

\usepackage[a4paper,margin=1.8cm,top=1.8cm,bottom=1.8cm,headheight=16pt]{geometry}
\usepackage{graphicx}
\usepackage{array}
\usepackage{booktabs}
\usepackage{tabularx}
\usepackage{multirow}
\usepackage{enumitem}
\usepackage{float}
\usepackage{titlesec}
\usepackage{fancyhdr}
\usepackage{setspace}
\usepackage{xcolor}
\usepackage{tikz}
\usepackage{tcolorbox}
\tcbuselibrary{skins}

% ---- Arabic/RTL packages (this order, before hyperref) ----
\usepackage{bidi}
\usepackage[no-math]{fontspec}
\usepackage[quiet]{polyglossia}

\setmainfont[Scale=1]{NotoKufiArabic-Regular.ttf}
\newfontfamily\arabicfont[Script=Arabic,Scale=0.95]{NotoKufiArabic-Regular.ttf}
\newfontfamily\englishfont[Scale=0.95]{TeX Gyre Termes}

\setdefaultlanguage[calendar=gregorian,hijricorrection=1]{arabic}
\setotherlanguage[variant=british]{english}

\usepackage[breaklinks,hidelinks]{hyperref}
\hypersetup{pdftitle={صفحة واحدة},pdfauthor={اسم الشركة}}

% ============================================================
% Accent color
% ============================================================
\definecolor{accent}{HTML}{6A1B9A}
\definecolor{accentdark}{HTML}{4A148C}
\definecolor{accentlight}{HTML}{F3E5F5}
\definecolor{statbg}{HTML}{4A148C}
\definecolor{panelbg}{HTML}{FAFAFA}
\definecolor{ctabg}{HTML}{6A1B9A}

% ============================================================
% Section formatting
% ============================================================
\titleformat{\section}{\normalsize\bfseries\color{accentdark}}{\thesection}{0.5em}{}
\titlespacing*{\section}{0pt}{0.6ex plus 0.1ex}{0.4ex plus 0.1ex}

% ============================================================
% Header / footer
% ============================================================
\pagestyle{fancy}
\fancyhf{}
\fancyhead[R]{\small\itshape صفحة واحدة (\LR{One-Pager})}
\fancyhead[L]{\small اسم الشركة}
\fancyfoot[C]{\small\thepage}
\renewcommand{\headrulewidth}{0pt}

% ============================================================
% Custom tcolorbox environments
% ============================================================
% Problem panel
\newtcolorbox{problembox}{
  enhanced,
  colback=accentlight,
  colframe=accent,
  boxrule=0.6pt,
  arc=3pt,
  left=10pt,right=10pt,top=6pt,bottom=6pt,
  title={\bfseries المشكلة},
  colbacktitle=accent,
  coltitle=white,
  fonttitle=\bfseries,
}

% Solution panel
\newtcolorbox{solutionbox}{
  enhanced,
  colback=panelbg,
  colframe=accentdark,
  boxrule=0.6pt,
  arc=3pt,
  left=10pt,right=10pt,top=6pt,bottom=6pt,
  title={\bfseries الحل},
  colbacktitle=accentdark,
  coltitle=white,
  fonttitle=\bfseries,
}

% Feature box
\newtcolorbox{featurebox}[1]{
  enhanced,
  colback=white,
  colframe=accent,
  boxrule=0.5pt,
  arc=2pt,
  left=6pt,right=6pt,top=4pt,bottom=4pt,
  title={#1},
  colbacktitle=accent,
  coltitle=white,
  fonttitle=\bfseries\small,
  attach title to upper,
  after title={\par\vspace{1pt}},
}

% Stat box for metrics row
\newtcolorbox{statbox}{
  enhanced,
  colback=statbg,
  colframe=statbg,
  boxrule=0pt,
  arc=4pt,
  left=4pt,right=4pt,top=6pt,bottom=6pt,
  coltext=white,
  halign=center,
}

% CTA footer box
\newtcolorbox{ctabox}{
  enhanced,
  colback=ctabg,
  colframe=ctabg,
  boxrule=0pt,
  arc=4pt,
  left=12pt,right=12pt,top=8pt,bottom=8pt,
  coltext=white,
}

\setstretch{1.15}
\setlength{\parindent}{0pt}
\setlength{\parskip}{0.2em}

\begin{document}
\thispagestyle{fancy}

% ============================================================
% HEADER (company/product name + tagline)
% ============================================================
\noindent\colorbox{accentdark}{\parbox{\dimexpr\textwidth-2\fboxsep\relax}{%
\vspace{0.3em}
\begin{center}
{\color{white}\Huge\bfseries اسم المنتج}\\[0.1em]
{\color{white}\large\itshape شعار أو وصف مختصر في جملة واحدة}
\end{center}
\vspace{0.2em}
}}

\vspace{0.5em}

% ============================================================
% PROBLEM & SOLUTION (side by side)
% ============================================================
\noindent
\begin{minipage}[t]{0.48\textwidth}
% TODO: صف المشكلة التي يحلها منتجك.
\begin{problembox}
وصف موجز للمشكلة التي يواجهها العملاء المستهدفون.%
اكتب جملتين أو ثلاثاً توضحان حجم المشكلة وأثرها.
\end{problembox}
\end{minipage}\hfill
\begin{minipage}[t]{0.48\textwidth}
% TODO: صف الحل الذي يقدمه منتجك.
\begin{solutionbox}
وصف موجز للحل الذي يقدمه اسم المنتج وكيف يعالج%
المشكلة المذكورة. اكتب جملتين أو ثلاثاً.
\end{solutionbox}
\end{minipage}

\vspace{0.4em}

% ============================================================
% KEY FEATURES
% ============================================================
\section{الميزات الرئيسية}

% TODO: استبدل بميزات منتجك (3--4 ميزات).
\noindent
\begin{minipage}[t]{0.23\textwidth}
\begin{featurebox}{\LR{[Icon]} الميزة 1}
وصف مختصر للميزة الأولى وفائدتها.
\end{featurebox}
\end{minipage}\hfill
\begin{minipage}[t]{0.23\textwidth}
\begin{featurebox}{\LR{[Icon]} الميزة 2}
وصف مختصر للميزة الثانية وفائدتها.
\end{featurebox}
\end{minipage}\hfill
\begin{minipage}[t]{0.23\textwidth}
\begin{featurebox}{\LR{[Icon]} الميزة 3}
وصف مختصر للميزة الثالثة وفائدتها.
\end{featurebox}
\end{minipage}\hfill
\begin{minipage}[t]{0.23\textwidth}
\begin{featurebox}{\LR{[Icon]} الميزة 4}
وصف مختصر للميزة الرابعة وفائدتها.
\end{featurebox}
\end{minipage}

\vspace{0.5em}

% ============================================================
% METRICS / STATS ROW
% ============================================================
% TODO: استبدل الأرقام بمؤشراتك الفعلية.
\noindent
\begin{minipage}[t]{0.31\textwidth}
\begin{statbox}
\centering
{\Huge\bfseries +10K}\\[0.1em]
{\small مستخدم نشط}
\end{statbox}
\end{minipage}\hfill
\begin{minipage}[t]{0.31\textwidth}
\begin{statbox}
\centering
{\Huge\bfseries 99\%}\\[0.1em]
{\small رضا العملاء}
\end{statbox}
\end{minipage}\hfill
\begin{minipage}[t]{0.31\textwidth}
\begin{statbox}
\centering
{\Huge\bfseries 24/7}\\[0.1em]
{\small دعم فني}
\end{statbox}
\end{minipage}

\vspace{0.5em}

% ============================================================
% TARGET MARKET
% ============================================================
\section{السوق المستهدف}
% TODO: صف السوق المستهدف.
اسم المنتج يستهدف ---وصف الشريحة المستهدفة--- في%
---القطاع أو المنطقة---. يبلغ حجم السوق المستهدف%
---رقم--- عميلاً محتملاً، مع معدل نمو سنوي يقدر بـ%
---رقم---\%.

\vspace{0.4em}

% ============================================================
% CONTACT / CTA FOOTER
% ============================================================
\begin{ctabox}
\begin{center}
{\Large\bfseries ابدأ اليوم!}\\[0.3em]
{\small جرّب اسم المنتج مجاناً لمدة 14 يوماً}
\end{center}
\vspace{0.3em}
\begin{center}
\LR{www.example.com} \quad|\quad%
\LR{\href{mailto:info@example.com}{info@example.com}} \quad|\quad%
\LR{+00-000-000-0000}
\end{center}
\end{ctabox}

\end{document}
